% corpus_supplement: Erdős problem #488
% source: https://www.erdosproblems.com/latex/488
% fetched_at: 2026-07-14T18:48:18Z
% payload_sha256: b9b6c006b2c87f8c8742a8686cbe8a3ea8c1f5dc592dafd228c3296727a47d7e
% statement/remarks/references extracted by erdos_ingest.extract_source_page;
% rendered by erdos_ingest.render_tex — do not edit
\documentclass{article}
\usepackage{amsmath, amssymb, amsthm}
\usepackage{hyperref}
\usepackage{geometry}
\geometry{margin=1in}

\title{Erd\H{o}s Problem \#488}
\date{Last updated: 2026-03-29}
\author{Source: erdosproblems.com}

\begin{document}
\maketitle

\noindent\textbf{Status:} OPEN \quad \textbf{No prize}

\noindent\textbf{Tags:} number theory

\medskip
\noindent\textbf{Problem Statement:}

Let $A$ be a finite set and\[B=\{ n \geq 1 : a\mid n\textrm{ for some }a\in A\}.\]Is it true that, for every $m>n\geq \max(A)$,\[\frac{\lvert B\cap [1,m]\rvert }{m}< 2\frac{\lvert B\cap [1,n]\rvert}{n}?\]

\bigskip
\noindent\textbf{Remarks:}

The constant $2$ would be the best possible here, as witnessed by taking $A=\{a\}$, $n=2a-1$, and $m=2a$.

This problem is also discussed in problem E5 of Guy's collection \cite{Gu04}.

In \cite{Er61} this problem is as stated above, but with $a\mid n$ in the definition of $B$ replaced by $a\nmid n$. This is most likely a typo (especially since the problem is also given as stated above in \cite{Er66}). There have been several counterexamples given for this alternate problem. Cambie has observed that, if $A$ is the set of primes bounded above by $n$, and $m=2n$, then\[\frac{\lvert B\cap [1,m]\rvert }{m}=\frac{\pi(2n)-\pi(n)+1}{2n}\sim \frac{1}{2\log n}\]while\[\frac{\lvert B\cap [1,n]\rvert}{n}=\frac{1}{n}.\]Further concrete counterexamples, found by Alexeev and Aristotle, are given in the comments section.

\bigskip
\noindent\textbf{References:}

[Er61] Erd\H{o}s, Paul, Some unsolved problems. Magyar Tud. Akad. Mat. Kutat\'{o} Int. K\"{o}zl. (1961), 221-254.
 
 [Er66] Erd\H{o}s, P\'{a}l, Remarks on number theory. {V}. {E}xtremal problems in number
theory. {II}. Mat. Lapok (1966), 135--155.
 
 [Gu04] Guy, Richard K., Unsolved problems in number theory. (2004), xviii+437.

\bigskip
\noindent\small{Source: \url{https://www.erdosproblems.com/488}}
\end{document}
